\newif\ifglobal

%%%%%%%%%%%%%%%%%%%%%%%
%%% Compile to view %%%
%%%%%%%%%%%%%%%%%%%%%%%

%%%%%%%%%%%%%%%%%%%%%%%%%%%%%%%%%%%%%%%%%%%%%%%%%%%
%%% Readme:                                     %%%
%%% https://www.overleaf.com/read/bwdjvfjksvjy/ %%%
%%%%%%%%%%%%%%%%%%%%%%%%%%%%%%%%%%%%%%%%%%%%%%%%%%%

% >>> M?: marks a module setting number or important notice

% >>> M4: Set {./relative-path} to external files for this module <<<
\def \modulefiles {./48-I2S}
% To add an external file, use:
% (Replace '---' with actual filename)
% '\input{\modulefiles/tables/---.tex}' for tables
% '\includegraphics{\modulefiles/figures/---}' for figures
% '\subfile{\modulefiles/---}' for register files

% >>> M1: This module is in global project? <<<
% 'YES' - uncomment; 'NO' - comment out
\globaltrue

\ifglobal
    \documentclass[main\_\_CN.tex]{subfiles}
   
\else
    % Font "FandolSong-Regular" used by ctex does not have GSUB table.
    % It causes warnings that do not affect compilation and can be ignored.
    \PassOptionsToPackage{quiet}{fontspec}

    \documentclass[openany, 10pt]{book}
   
    % >>> M2: In ./chip-spec-settings/preamble-trm-module, also find
    %         settings for visibility of todo notes, watermark <<< 
    \usepackage{preamble-trm-module}

    % Enable Chinese typesetting
    \usepackage{ctex}

    % >>> M3: Temporary labels in Chinese
    %         you can add your own labels there <<<
    \myexternaldocument{temp-labels-cn}

\fi %\ifglobal

\setcounter{secnumdepth}{4}

% >>> M5: If this module needs any updates in
% './chip-spec-settings/preamble-trm-module.sty'
% (include additional package, etc.), contact your module PM
% or create the file './chip-spec-settings/changelog.tex' make the
% necessary updates yourself and document them in this file <<<

\raggedright

\begin{document}

% Sets language into which pre-defined bits of text are translated
\selectlanguage{chinese}

% Do not delete this block
\ifglobal
% Add the following front matter if
% this module is outside of global project
\else
    \InputIfFileExists{readme}{}{}
    {\let\clearpage\relax \listoftodos}
    {\let\clearpage\relax \listofdonetodos}
    \tableofcontents
    %\listoftables
    %\listoffigures
\fi


%%%%%%%%%%%%%%%%%%%%%%%%%
%%% TRM Module Begins %%%
%%%%%%%%%%%%%%%%%%%%%%%%%

\hypertarget{i2s}{}
\chapter{I2S 控制器 (I2S)}
\label{mod:i2s}

\section{概述}

\chipname{} 内置三个 I2S 接口，即 I2S\regindex{0}、I2S\regindex{1} 和 I2S\regindex{2}，为多媒体应用，尤其是为数字音频应用提供了灵活的数据通信接口。

I2S 标准总线定义了三种信号——串行时钟信号 BCK、字选择信号 WS 和串行数据信号 SD。一个基础的 I2S 数据总线包含一个主机和一个从机。主机和从机的角色在通信过程中保持不变。\chipname{} 的 I2S 模块包含独立的 TX 模块和 RX 模块，能够保证优良的通信性能。

\begin{tiplisting}
本章节提供的信息适用于 I2S\regindex{0}、I2S\regindex{1} 和 I2S\regindex{2}。除非另有说明，本章节中 I2S 及 I2S\regindex{n} 均指代 I2S\regindex{0}、I2S\regindex{1} 和 I2S\regindex{2}。
\end{tiplisting}
\vspace{-2em}

\section{术语}
\label{sec:i2s-term}

为了更好地说明 I2S 的功能，本章使用了以下术语。 

\begin{longtable}[c]{ R{4cm} p{10cm} }

\textbf{主机模式}     & I2S\regindex{n} 作为主机驱动 BCK/WS，向从机发送或从其接收数据。 \\ 
\textbf{从机模式}      & I2S\regindex{n} 作为从机由 BCK/WS 驱动，从主机接收或向其发送数据。 \\
\textbf{全双工} & 主机与从机之间的发送线和接收线各自独立，发送数据和接收数据同时进行。\\
\textbf{半双工} & 主机和从机只能有一方先发送数据，另一方接收数据。发送数据和接收数据不能同时进行。 \\
\textbf{A 律和 $\mu$ 律} & A 律和 \textrm{$\mu$} 律是数字脉冲编码调制 (PCM) 非均匀量化中的压缩/解压缩算法，可以有效地改善信号的量化信噪比。 \\
\textbf{TDM RX 模式} & 即 TDM 输入模式，该模式利用时分复用方式接收脉冲编码调制 (PCM) 数据。信号线包括 BCK、WS 和 DATA。可以接收最多 16 个通道的数据。用户可配置数据格式，支持 TDM Philips 格式、TDM MSB 对齐格式、TDM PCM 格式等。\\
\textbf{PDM RX 模式} & 即 PDM 输入模式，该模式接收脉冲密度调制 (PDM) 数据。信号线包括 WS 和 DATA。用户可配置数据格式，支持 PDM 标准格式等。\\
\textbf{TDM TX 模式} & 即 TDM 输出模式，该模式利用时分复用方式发送脉冲编码调制 (PCM) 数据。信号线包括 BCK、WS 和 DATA，可以发送最多 16 个通道的数据。用户可配置数据格式，支持 TDM Philips 格式、TDM MSB 对齐格式、TDM PCM 格式等。\\
\textbf{PDM TX 模式} & 即 PDM 输出模式，该模式发送脉冲密度调制 (PDM) 数据。信号线包括 WS 和 DATA。用户可配置数据格式，支持 PDM 标准格式等。\\
\textbf{PCM-to-PDM 数据格式转换器（仅适用于 I2S\regindex{0}）} & 将 PCM 数据转换为 PDM 数据的格式转换器（后文简称为 PCM-to-PDM 转换器），可以在 PDM TX 主机模式下启用。启用后，I2S 通过 GDMA 从存储器中取得脉冲编码调制 (PCM) 数据，将其转换为脉冲密度调制 (PDM) 数据并发送。\\
\textbf{PDM-to-PCM 数据格式转换器（仅适用于 I2S\regindex{0}）} & 将 PDM 数据转换为 PCM 数据的格式转换器（后文简称为 PDM-to-PCM 转换器），可以在 PDM RX 主机或从机模式下启用。启用后，I2S 接收脉冲密度调制 (PDM) 数据，将其转换为脉冲编码调制 (PCM) 数据，并将其通过 GDMA 存入存储器。\\
\end{longtable}

\section{特性}

I2S 模块具有以下特性：

\begin{itemize}
    \item 支持主机模式和从机模式
    \item 支持全双工和半双工通信
    \item 支持 TX 模块和 RX 模块独立工作或同时工作
    \item 支持多种音频标准：
        \begin{itemize}    
            \item TDM Philips 标准
            \item TDM MSB 对齐标准
            \item TDM PCM 标准
            \item PDM 标准
        \end{itemize}
    \item 支持多种 TX/RX 模式
        \begin{itemize}
        \item TDM TX 模式，最多支持 16 通道
        \item TDM RX 模式，最多支持 16 通道
        \item PDM TX 模式
        \begin{itemize}
            \item 支持原始 PDM 数据发送
            \item 支持将 PCM 数据转换为 PDM 数据发送（仅适用于 I2S\regindex{0}），最多支持 2 通道
        \end{itemize}
        \item PDM RX 模式
        \begin{itemize}
            \item 支持原始 PDM 数据接收
            \item 支持将 PDM 数据转换为 PCM 数据接收（仅适用于 I2S\regindex{0}），最多支持 8 通道
        \end{itemize}
        \end{itemize}
    \item 可配置 APLL 时钟源，支持最高频率为 125 MHz
    \item 可配置高精度采样时钟，支持多种采样频率（详见下文注释）
    \item 支持 8/16/24/32 位的数据位宽
    \item TX 模式支持同步计数器
    \item 支持 ETM 功能
    \item 支持 GDMA（仅 GDMA-AHB）
    \item 支持 I2S 接口中断
\end{itemize}

\begin{tiplisting}
从机模式下，为了保障采样的准确性，模块时钟频率必须大于等于 BCK 时钟频率的 8 倍，因此 I2S\regindex{n} 支持的采样频率受到数据位宽和声道数的限制。例如，时钟源频率最高为 125 MHz，即模块时钟最高能配置为 62.5 MHz，能支持的 BCK 时钟频率最高为 7.8125 MHz，因此传输双声道 32 位宽数据时，能够支持 122.070 kHz 以下的采样频率，例如 8 kHz、16 kHz、32 kHz、44.1 kHz、48 kHz、88.2 kHz、96 kHz。具体请参考章节 \ref{sec:i2s-tx-rx-clk} \textit{\nameref{sec:i2s-tx-rx-clk}}。
\end{tiplisting}

\section{系统架构}

图 \ref{fig:i2s-sys-diagram} 为 \chipname{} I2S\regindex{n} 模块的结构框图。\chipname{} I2S\regindex{n} 模块包含：

\begin{itemize}
    \item 独立的 TX 模块
    \item 独立的 RX 模块
    \item 输入输出时序调节单元（I/O 同步）
    \item 时钟分频器（时钟发生器）
    \item 64 x 32 位 TX FIFO
    \item 64 x 32 位 RX FIFO
    \item 压缩/解压缩模块
\end{itemize}

\iftagged{ESP32-P4-latest}{%
\begin{figure}[H]
    \centering
    \fbox{\includegraphics[width=0.9\textwidth]{\modulefiles/figures/i2s_arch_20250808__CN}}
    \\ 注：图中 PDM-to-PCM 转换器和 PCM-to-PDM 转换器仅适用于 I2S\regindex{0}。
    \caption{\chipname{} I2S 系统框图}
    \label{fig:i2s-sys-diagram}
\end{figure}
}{%
\begin{figure}[H]
    \centering
    \fbox{\includegraphics[width=0.9\textwidth]{\modulefiles/figures/i2s_arch_20240201__CN}}
    \\ 注：图中 PDM-to-PCM 转换器和 PCM-to-PDM 转换器仅适用于 I2S\regindex{0}。
    \caption{\chipname{} I2S 系统框图}
    \label{fig:i2s-sys-diagram}
\end{figure}
}%

\chipname{} I2S\regindex{n} 模块支持 GDMA，可访问内部存储器，更多信息见章节 \ref{mod:dma} \textit{\nameref{mod:dma}}。

TX 模块和 RX 模块各自有一组三线接口，分别为串行时钟线 BCK、字选择线 WS 和串行数据线 SD。其中，TX 模块的 SD 线固定为输出，RX 模块的 SD 线固定为接收。主机模式下，TX 模块和 RX 模块的 BCK 和 WS 信号线为输出，从机模式下，TX 模块和 RX 模块的 BCK 和 WS 信号线为输入。

图 \ref{fig:i2s-sys-diagram} 中的右半部分展示了 I2S\regindex{n} 模块的信号总线。RX 和 TX 模块的信号命名规则为：I2S\regindex{n}\textcolor{red}{A}\_\textcolor{red}{B}\_\textcolor{red}{C}，例如 I2S\regindex{n}\textcolor{red}{I}\_\textcolor{red}{BCK}\_\textcolor{red}{in}。其中：

\begin{itemize}
    \item ``A" 表示信号属于 I2S\regindex{n} TX 或 RX 模块，包括：
    \begin{itemize}
        \item ``I" 表示信号输入至/输出自 RX 模块
        \item ``O" 表示信号输入至/输出自 TX 模块
    \end{itemize}
    \item ``B" 表示信号功能，包括：
    \begin{itemize}
        \item BCK
        \item WS
        \item SD
    \end{itemize}
    \item ``C" 表示该信号的方向，包括：
    \begin{itemize}
        \item ``in" 表示该信号输入 I2S\regindex{n} 模块
        \item ``out" 表示该信号自 I2S\regindex{n} 模块输出
    \end{itemize}
\end{itemize}

I2S 各信号的具体描述见表 \ref{tab:i2s-signal-descr}。

\begin{table}[H]
    \centering
    \caption{模块信号描述}
    \label{tab:i2s-signal-descr}
    \begin{threeparttable}
    \begin{tabular}{|p{3cm}|p{1.5cm}|p{8cm}| }
    \hline
    \rowcolor{lightgray}
    \textbf{信号}&\textbf{方向}&\textbf{功能}\\
    \hline
    I2S\regindex{n}I\_BCK\_in & 输入 & I2S\regindex{n} 从机模式下，输入 BCK 信号，用于 RX 模块\\
    \hline
    I2S\regindex{n}I\_BCK\_out& 输出 & I2S\regindex{n} 主机模式下，输出 BCK 信号，用于 RX 模块\\
    \hline
    I2S\regindex{n}I\_WS\_in & 输入 & I2S\regindex{n} 从机模式下，输入 WS 信号，用于 RX 模块\\
    \hline
    I2S\regindex{n}I\_WS\_out& 输出 & I2S\regindex{n} 主机模式下，输出 WS 信号，用于 RX 模块\\\hline
    I2S\regindex{n}I\_Data\_in& 输入 & I2S\regindex{n} RX 模块 的串行输入数据线\\\hline
    I2S\regindex{n}O\_Data\_out& 输出 & I2S\regindex{n} TX 模块 的串行输出数据线\\\hline
    I2S\regindex{n}O\_BCK\_in& 输入 & I2S\regindex{n} 从机模式下，输入 BCK 信号，用于 TX 模块 \\\hline
    I2S\regindex{n}O\_BCK\_out& 输出 & I2S\regindex{n} 主机模式下，输出 BCK 信号，用于 TX 模块 \\\hline
    I2S\regindex{n}O\_WS\_in& 输入 & I2S\regindex{n} 从机模式下，输入 WS 信号，用于 TX 模块\\\hline
    I2S\regindex{n}O\_WS\_out& 输出 & I2S\regindex{n} 主机模式下，输出 WS 信号，用于 TX 模块\\\hline
    I2S\regindex{n}\_MCLK\_in& 输入 & I2S\regindex{n} 从机模式下，来自外部芯片的时钟源\\\hline
    I2S\regindex{n}\_MCLK\_out& 输出 & I2S\regindex{n} 主机模式下，作为外部芯片的时钟源\\\hline
    I2S\regindex{0}I\_Data1\_in& 输入 & RX 模块在启用 PDM-to-PCM 转换器时的串行输入数据线\\\hline
    I2S\regindex{0}I\_Data2\_in& 输入 & RX 模块在启用 PDM-to-PCM 转换器时的串行输入数据线\\\hline
    I2S\regindex{0}I\_Data3\_in& 输入 & RX 模块在启用 PDM-to-PCM 转换器时的串行输入数据线\\\hline
    I2S\regindex{0}O\_Data1\_out& 输出 & TX 模块在启用 PCM-to-PDM 转换器时的串行输出数据线\\\hline
    \end{tabular}
    \begin{tablenotes}
            \item[*] I2S\regindex{n} 的所有信号均需要经过 HP GPIO 交换矩阵映射到芯片的管脚。更多信息请参考章节 \ref{mod:iomuxgpio} \textit{\nameref{mod:iomuxgpio}}。
    \end{tablenotes}
    \end{threeparttable}
\end{table}


\section{I2S 模块支持的音频协议} \label{sec:i2s-supported-audio-standards}

\chipname{} I2S\regindex{n} 模块支持多种音频标准，包括 TDM Philips 标准、TDM MSB 对齐标准、TDM PCM 标准以及 PDM 标准。

用户可通过配置以下寄存器，选择所需的音频标准：

\begin{itemize}
    \item \hyperref[fielddesc:I2SRXMSBSHIFT]{I2S\_TX/RX\_TDM\_EN}
    \begin{itemize}
        \item 0：禁用 TDM 模式
        \item 1：选择 TDM 模式
    \end{itemize}
    \item \hyperref[fielddesc:I2SRXMSBSHIFT]{I2S\_TX/RX\_PDM\_EN}
    \begin{itemize}
        \item 0：禁用 PDM 模式
        \item 1：选择 PDM 模式
    \end{itemize}    
    \item \hyperref[fielddesc:I2SRXMSBSHIFT]{I2S\_TX/RX\_MSB\_SHIFT}
    \begin{itemize}
        \item 0：配置 WS 信号和 SD 信号同时开始变化，即选择 MSB 对齐标准
        \item 1：配置 WS 信号先于 SD 信号一个 BCK 时钟周期开始变化，即选择 Philips 标准或 PCM 标准
    \end{itemize}
\end{itemize}

\begin{tiplisting}
不可将 \hyperref[fielddesc:I2SRXMSBSHIFT]{I2S\_TX/RX\_TDM\_EN} 和 \hyperref[fielddesc:I2SRXMSBSHIFT]{I2S\_TX/RX\_PDM\_EN} 同时配置为 1 或同时配置为 0，否则会导致 \chipname{} I2S\regindex{n} 以既非 TDM 模式也非 PDM 模式的状态错误发送数据。
\end{tiplisting}

\subsection{TDM Philips 标准模式}

在 Philips 标准下，在 BCK 的下降沿，WS 信号先于 SD 信号一个 BCK 时钟周期开始变化，即 WS 信号从当前通道数据的第一个位之前的一个时钟开始有效，并在当前通道数据发送结束前一个 BCK 时钟周期开始变化。SD 信号线上首先传输音频数据的最高位。

与 Philips 标准相比，TDM Philips 标准支持的通道更多，见图 \ref{fig:i2s-tdm-philips-std-timing-diagram}。

\begin{figure}[H]
    \centering
    \includegraphics[width=1.0\textwidth]{\modulefiles/figures/i2s_tdm_phillips_mode.png}
    \caption{时序图 -- TDM Philips 标准}
    \label{fig:i2s-tdm-philips-std-timing-diagram}
\end{figure}

\subsection{TDM MSB 对齐标准模式}

MSB 对齐标准下，在 BCK 下降沿，WS 信号和 SD 信号同时变化。WS 持续到当前通道数据发送结束，SD 信号线上首先传输音频数据的最高位。

与 MSB 对齐标准相比，TDM MSB 对齐标准支持更多通道，见图 \ref{fig:i2s-tdm-msb-align-std-timing-diagram}。

\begin{figure}[H]
    \centering
    \includegraphics[width=1.0\textwidth]{\modulefiles/figures/i2s_tdm_msb_mode.png}
    \caption{时序图 -- TDM MSB 对齐标准}
    \label{fig:i2s-tdm-msb-align-std-timing-diagram}
\end{figure}

\subsection{TDM PCM 标准模式}

在 PCM 标准的短帧同步模式下，在 BCK 的下降沿，WS 信号先于 SD 信号一个 BCK 时钟周期开始变化，即 WS 信号从当前通道数据的第一个位之前的一个时钟开始有效，并持续一个 BCK 时钟周期。SD 信号线上首先传输音频数据的最高位。

与 PCM 标准相比，TDM PCM 标准支持更多通道，见图 \ref{fig:i2s-tdm-pcm-std-timing-diagram} 所示。

\begin{figure}[H]
    \centering
    \includegraphics[width=1.0\textwidth]{\modulefiles/figures/i2s_tdm_pcm_mode.png}
    \caption{时序图 -- TDM PCM 标准}
    \label{fig:i2s-tdm-pcm-std-timing-diagram}
\end{figure}


\subsection{PDM 标准模式}

如图 \ref{fig:i2s-pdm-std-timing-diagram} 所示，在 PDM 标准下，WS 代表左/右声道，WS 与 SD 同时变化。WS 在数据发送过程中持续变化，WS 的高低对应两个声道。请参考图 \ref{fig:i2s-pdm-std-timing-diagram}。

\begin{figure}[H]
    \centering
    \includegraphics[width=0.8\textwidth]{\modulefiles/figures/i2s_pdm_mode.pdf}
    \caption{时序图 -- PDM 标准}
    \label{fig:i2s-pdm-std-timing-diagram}
\end{figure}


\section{I2S TX/RX 模块时钟}\label{sec:i2s-tx-rx-clk}

I2S\regindex{n}\_TX/RX\_CLK 作为 I2S\regindex{n} TX/RX 模块的主时钟，可由以下时钟分频所得：

\begin{itemize}
    \item 40 MHz XTAL\_CLK
    \iftagged{ESP32-P4-latest}{\item 160 MHz PLL\_F160M\_CLK}{}%
    \item 可调频率的 APLL\_CLK
    \item 外部输入时钟 I2S\regindex{n}\_MCLK\_in 
\end{itemize}

I2S\regindex{n} TX/RX 模块的串行时钟 BCK 再由 I2S\regindex{n}\_TX/RX\_CLK 分频获得，如图 \ref{fig:i2s-clk} 所示。\hyperref[fielddesc:HPSYSCLKRSTI2S0RXCLKSRCSEL]{HP\_SYS\_CLKRST\_I2S\regindex{n}\_TX/RX\_CLK\_SRC\_SEL} 用于选择 I2S TX/RX 的时钟源，\hyperref[fielddesc:HPSYSCLKRSTI2S0RXCLKEN]{HP\_SYS\_CLKRST\_I2S\regindex{n}\_TX/RX\_CLK\_EN} 用于使能或者关闭 I2S TX/RX 模块的时钟源。

\iftagged{ESP32-P4-latest}{%
\begin{figure}[H]
    \centering
    \fbox{\includegraphics[width=0.8\textwidth]{\modulefiles/figures/i2s_clk_20250808}}
    \caption{I2S\regindex{n} 时钟分频器}
    \label{fig:i2s-clk}
\end{figure}
}{%
\begin{figure}[H]
    \centering
    \fbox{\includegraphics[width=0.8\textwidth]{\modulefiles/figures/i2s_clk_20231128}}
    \caption{I2S\regindex{n} 时钟分频器}
    \label{fig:i2s-clk}
\end{figure}
}%

I2S\regindex{n}\_TX/RX\_CLK 的频率 $f_{\textrm{I2S\regindex{n}\_TX/RX\_CLK}}$ 与分频器时钟源频率 $f_{\textrm{I2S\regindex{n}\_CLK\_S}}$ 间的关系如下：
$$f_{\textrm{I2S\regindex{n}\_TX/RX\_CLK}}=\frac{f_{\textrm{I2S\regindex{n}\_CLK\_S}}}{\textrm{N} + \frac{\textrm{b}}{\textrm{a}}}$$

其中，2 <= N <= 256（N 为整数），N 对应寄存器域 \hyperref[fielddesc:HPSYSCLKRSTI2S0TXDIVN]{HP\_SYS\_CLKRST\_I2S\regindex{n}\_TX/RX\_DIV\_N} 的值，具体为：
\begin{itemize}
    \item \hyperref[fielddesc:HPSYSCLKRSTI2S0TXDIVN]{HP\_SYS\_CLKRST\_I2S\regindex{n}\_TX/RX\_DIV\_N} = 0 时，N = 256；
    \item \hyperref[fielddesc:HPSYSCLKRSTI2S0TXDIVN]{HP\_SYS\_CLKRST\_I2S\regindex{n}\_TX/RX\_DIV\_N} = 1 时，N = 2；
    \item \hyperref[fielddesc:HPSYSCLKRSTI2S0TXDIVN]{HP\_SYS\_CLKRST\_I2S\regindex{n}\_TX/RX\_DIV\_N} 为其他值时，N = \hyperref[fielddesc:HPSYSCLKRSTI2S0TXDIVN]{HP\_SYS\_CLKRST\_I2S\regindex{n}\_TX/RX\_DIV\_N} 的值。
\end{itemize}

分数部分分频系数 a、b 唯一对应系数 x、y、z、yn1。系数对应公式为：

\begin{itemize}
    \item 当 $\textrm{b} <= \frac{\textrm{a}}{2}$ 时，$\textrm{yn1} = 0$，$\textrm{x} = floor([\frac{\textrm{a}}{\textrm{b}}]) - 1$，$\textrm{y} = \textrm{a}\%\textrm{b}$，$\textrm{z} = \textrm{b}$；
    \item 当 $\textrm{b} > \frac{\textrm{a}}{2}$ 时，$\textrm{yn1}=1$，$\textrm{x} = floor([\frac{\textrm{a}}{\textrm{a - b}}]) - 1$，$\textrm{y} = \textrm{a}\%\textrm{(a - b)}$，$\textrm{z}=\textrm{a - b}$。
\end{itemize}

系数 x、y、z、yn1 在 \hyperref[fielddesc:HPSYSCLKRSTI2S0TXDIVX]{HP\_SYS\_CLKRST\_I2S\regindex{n}\_TX/RX\_DIV\_X}、\hyperref[fielddesc:HPSYSCLKRSTI2S0TXDIVY]{HP\_SYS\_CLKRST\_I2S\regindex{n}\_TX/RX\_DIV\_Y}、\hyperref[fielddesc:HPSYSCLKRSTI2S0TXDIVZ]{HP\_SYS\_CLKRST\_I2S\regindex{n}\_TX/RX\_DIV\_Z}、\hyperref[fielddesc:HPSYSCLKRSTI2S0TXDIVYN1]{HP\_SYS\_CLKRST\_I2S\regindex{n}\_TX/RX\_DIV\_YN1} 中配置。 

对于整数分频，\hyperref[fielddesc:HPSYSCLKRSTI2S0TXDIVX]{HP\_SYS\_CLKRST\_I2S\regindex{n}\_TX/RX\_DIV\_X} 和 \hyperref[fielddesc:HPSYSCLKRSTI2S0TXDIVZ]{HP\_SYS\_CLKRST\_I2S\regindex{n}\_TX/RX\_DIV\_Z} 清零，\hyperref[fielddesc:HPSYSCLKRSTI2S0TXDIVY]{HP\_SYS\_CLKRST\_I2S\regindex{n}\_TX/RX\_DIV\_Y} 置为 1。

\begin{tiplisting}
使用小数分频功能可能会产生时钟抖动。
\end{tiplisting}

在主机发送模式下，I2S\regindex{n} TX 模块的串行时钟 BCK 为 I2S\regindex{n}O\_BCK\_out 信号，由 I2S\regindex{n}\_TX\_CLK 分频获得。即：

$$f_{\textrm{I2S\regindex{n}O\_BCK\_out}} = \frac{f_{\textrm{I2S\regindex{n}\_TX\_CLK}}}{\textrm{MO}}$$

其中，MO 的值为：
$$ \textrm{MO} = \textrm{\hyperref[fielddesc:I2STXBCKDIVNUM]{I2S\regindex{n}\_TX\_BCK\_DIV\_NUM}} + 1 $$

\vspace{-1em}
\begin{tiplisting}
注意，\hyperref[fielddesc:I2STXBCKDIVNUM]{I2S\_TX\_BCK\_DIV\_NUM} 不可配置为 1。
\end{tiplisting}

在主机接收模式下，I2S\regindex{n} RX 模块的串行时钟 BCK 为 I2S\regindex{n}I\_BCK\_out 信号，由 I2S\regindex{n}\_RX\_CLK 分频获得。即：

$$f_{\textrm{I2S\regindex{n}I\_BCK\_out}} = \frac{f_{\textrm{I2S\regindex{n}\_RX\_CLK}}}{\textrm{MI}}$$

其中 MI 的值为：
$$ \textrm{MI} = \textrm{\hyperref[fielddesc:I2STXBCKDIVNUM]{I2S\_RX\_BCK\_DIV\_NUM}} + 1 $$
\begin{tiplisting}
\vspace{-2em}
\begin{itemize}
    \item \hyperref[fielddesc:I2SRXBCKDIVNUM]{I2S\_RX\_BCK\_DIV\_NUM} 不可配置为 1；
    \item 当模块处于 I2S 从机模式时，必须保证 $f_{\textrm{I2S\_TX/RX\_CLK}}$ >= 8 × $f_{\textrm{BCK}}$。另外模块可以输出 I2S\regindex{n}\_MCLK\_out 作为外部设备的主时钟。
\end{itemize}
\end{tiplisting}

\section{I2S 模块复位} \label{sec:i2s-reset}

I2S\regindex{n} 模块中各个单元以及 FIFO 可通过配置相关位进行复位：

\begin{itemize}
\item I2S\regindex{n} TX/RX 单元：可配置 \hyperref[fielddesc:I2STXRESET]{I2S\_TX\_RESET} 和 \hyperref[fielddesc:I2SRXRESET]{I2S\_RX\_RESET} 进行复位；
\item I2S\regindex{n} TX/RX FIFO：可配置 \hyperref[fielddesc:I2STXFIFORESET]{I2S\_TX\_FIFO\_RESET} 和 \hyperref[fielddesc:I2SRXFIFORESET]{I2S\_RX\_FIFO\_RESET} 进行复位。
\end{itemize}

\begin{tiplisting}
在模块和 FIFO 复位之前，需要先配置 I2S\regindex{n} 模块时钟。
\end{tiplisting} 


\section{I2S 主/从机模式}

\chipname{} I2S\regindex{n} 模块可作为主机或从机，两种模式均支持半双工通信或全双工通信。用户可配置  \hyperref[fielddesc:I2SRXSLAVEMOD]{I2S\_RX\_SLAVE\_MOD} 和 \hyperref[fielddesc:I2STXSLAVEMOD]{I2S\_TX\_SLAVE\_MOD} 选择需要的模式。

\begin{itemize}
    \item \hyperref[fielddesc:I2STXSLAVEMOD]{I2S\_TX\_SLAVE\_MOD} 
    \begin{itemize}
        \item 0：主机发送模式
        \item 1：从机发送模式
    \end{itemize}
    \item \hyperref[fielddesc:I2SRXSLAVEMOD]{I2S\_RX\_SLAVE\_MOD}
    \begin{itemize}
        \item 0：主机接收模式
        \item 1：从机接收模式
    \end{itemize}
\end{itemize}

\subsection{主/从机发送模式} \label{sec:i2s-master-slave-tx-mode}
\begin{itemize}
    \item 主机发送模式
    \begin{itemize}
        \item 置位 \hyperref[fielddesc:I2STXSTART]{I2S\_TX\_START} 启动一次发送操作。置位该位，TX 模块会一直输出时钟信号和串行数据。
        \item 置位 \hyperref[fielddesc:I2STXSTOPEN]{I2S\_TX\_STOP\_EN} 时，如果 FIFO 中的数据全部发送完毕，则主机停止发送数据和时钟信号。
        \item 清零 \hyperref[fielddesc:I2STXSTOPEN]{I2S\_TX\_STOP\_EN} 时，如果 FIFO 中的数据全部发送完毕，并且没有新数据填入，发送模块将一直发送最后一帧数据和时钟信号。
        \item 当 \hyperref[fielddesc:I2STXSTART]{I2S\_TX\_START} 被清零时，主机停止发送数据。
    \end{itemize}
    \item 从机发送模式
     \begin{itemize}
        \item 置位 \hyperref[fielddesc:I2STXSTART]{I2S\_TX\_START}。TX 模块等待主机 BCK 时钟，来启动发送操作。
        \item 置位 \hyperref[fielddesc:I2STXSTOPEN]{I2S\_TX\_STOP\_EN} 时，如果 FIFO 中的数据全部发送完毕，则从机发送数据一直为零，直到主机停止发送 BCK 时钟为止。
        \item 清零 \hyperref[fielddesc:I2STXSTOPEN]{I2S\_TX\_STOP\_EN} 时，如果 FIFO 中的数据全部发送完毕，并且没有新数据填入， TX 模块将一直发送最后一帧数据。
        \item 当 \hyperref[fielddesc:I2STXSTART]{I2S\_TX\_START} 被清零时，从机发送数据一直为零，直到主机停止发送 BCK 时钟为止。
    \end{itemize}
\end{itemize}

\subsection{主/从机接收模式}
\begin{itemize}
    \item 主机接收模式
    \begin{itemize}
        \item 置位 \hyperref[fielddesc:I2SRXSTART]{I2S\_RX\_START} 启动一次接收操作。RX 模块会一直输出时钟信号，并对输入数据进行采样。
        \item 配置 \hyperref[fielddesc:I2SRXSTOPMODE]{I2S\_RX\_STOP\_MODE} 来控制接收模式何时暂停接收：
        \begin{itemize}
            \item 0：只有 \hyperref[fielddesc:I2SRXSTART]{I2S\_RX\_START} 被清除，RX 模块才会暂停接收。
            \item 1：RX 模块会在 \hyperref[fielddesc:I2SRXSTART]{I2S\_RX\_START} 被清除，或者接收字节数大于 \hyperref[regdesc:I2SRXEOFNUMREG]{I2S\_RX\_EOF\_NUM\_REG} 配置的接收长度值时暂停接收。当 \hyperref[fielddesc:I2SRXSTART]{I2S\_RX\_START} 未被清除时，可通过置位 \hyperref[fielddesc:I2SRXRESET]{I2S\_RX\_RESET} 来重新启动接收。
            \item 2：RX 模块会在 \hyperref[fielddesc:I2SRXSTART]{I2S\_RX\_START} 被清除或 GDMA RX FIFO 满时暂停接收。当 \hyperref[fielddesc:I2SRXSTART]{I2S\_RX\_START} 未被清除且 GDMA RX FIFO 不再满时，可通过置位 \hyperref[fielddesc:I2SRXRESET]{I2S\_RX\_RESET} 重新启动接收。
        \end{itemize}
        \item 清零 \hyperref[fielddesc:I2SRXSTART]{I2S\_RX\_START}，RX 模块停止接收数据。
    \end{itemize}
    \item 从机接收模式
    \begin{itemize}
        \item 置位 \hyperref[fielddesc:I2SRXSTART]{I2S\_RX\_START}。等待主机 BCK 时钟，来启动接收操作。
        \item 配置 \hyperref[fielddesc:I2SRXSTOPMODE]{I2S\_RX\_STOP\_MODE} 来控制接收模式何时暂停接收：
        \begin{itemize}
            \item 0：只有 \hyperref[fielddesc:I2SRXSTART]{I2S\_RX\_START} 被清除，RX 模块才会暂停接收；
            \item 1：RX 模块会在 \hyperref[fielddesc:I2SRXSTART]{I2S\_RX\_START} 被清除，或者接收字节数大于 \hyperref[regdesc:I2SRXEOFNUMREG]{I2S\_RX\_EOF\_NUM\_REG} 配置的接收长度值时暂停接收。当 \hyperref[fielddesc:I2SRXSTART]{I2S\_RX\_START} 未被清除时，可通过置位 \hyperref[fielddesc:I2SRXRESET]{I2S\_RX\_RESET} 来重新开始接收；
            \item 2：RX 模块会在 \hyperref[fielddesc:I2SRXSTART]{I2S\_RX\_START} 被清除或 GDMA RX FIFO 满时暂停接收。当 \hyperref[fielddesc:I2SRXSTART]{I2S\_RX\_START} 未被清除且 GDMA RX FIFO 不再满时，可通过置位 \hyperref[fielddesc:I2SRXRESET]{I2S\_RX\_RESET} 恢复接收。
        \end{itemize}
        \item 清零 \hyperref[fielddesc:I2SRXSTART]{I2S\_RX\_START}，RX 模块停止接收数据。
    \end{itemize}
\end{itemize}

\section{发送数据}\label{sec:i2s-tx-data}

\vspace{1em}
\begin{tiplisting}
本小节以及后续小节所述的配置，均需要通过置位 \hyperref[fielddesc:I2STXUPDATE]{I2S\_TX\_UPDATE} 的方式来进行更新，从而将 I2S\regindex{n} TX 寄存器数据从 APB 时钟域同步到 I2S\regindex{n} TX 时钟域。详细配置见第 \ref{sec:i2s-configure-i2s-tx-mode} 小节。
\end{tiplisting}
\vspace{-1em}

\chipname{} I2S\regindex{n} 发送数据时，从 GDMA 读取数据，经过数据格式控制和通道模式控制，从外设输出信号输出对应数据。
\subsection{数据格式控制}\label{sec:i2s-data-control}
数据格式控制分为三个阶段：
\begin{itemize}
    \item 第一阶段：从内存中读出有效数据并写入 TX FIFO；
    \item 第二阶段：将待发送数据从 TX FIFO 中读出，并进行输出数据模式转换；
    \item 第三阶段：将待发送数据转换为串行数据流输出。
\end{itemize}

\subsubsection{通道有效数据位宽}
\hyperref[fielddesc:I2STXBITSMOD]{I2S\_TX\_BITS\_MOD} 和 \hyperref[fielddesc:I2STX24FILLEN]{I2S\_TX\_24\_FILL\_EN} 决定了每个通道的有效数据位宽，其可取值和对应的有效数据位宽如下表：

\begin{table}[H]
    \centering
    \caption{通道有效数据位宽控制}
    \label{tab:i2s-bit-width-chan-valid-data}
    \begin{threeparttable}
    \begin{tabular}{|p{4.2cm}|R{3.1cm}|R{3.5cm}|}
    \hline
    \rowcolor{lightgray}
    \textbf{通道有效数据位宽} & \textbf{\hyperref[fielddesc:I2STXBITSMOD]{I2S\_TX\_BITS\_MOD}}  & \textbf{\hyperref[fielddesc:I2STX24FILLEN]{I2S\_TX\_24\_FILL\_EN}} \\ \hline
    \multirow{2}*{32} & 31 & x\tnote{*} \\\cline{2-3}
                      & 23 & 1 \\\hline

                24    & 23 & 0 \\\hline

                16     & 15 & x \\\hline
                8     & 7 & x \\\hline
    \end{tabular}
            \begin{tablenotes}
            \item[*] ``x" 表示该值被忽略。
        \end{tablenotes}

    \end{threeparttable}
\end{table}

\subsubsection{通道有效数据字节序}
I2S\regindex{n} 通过 GDMA 读取数据之后，\hyperref[fielddesc:I2STXBIGENDIAN]{I2S\_TX\_BIG\_ENDIAN} 用于控制从 GDMA 读取数据的字节序。表 \ref{tab:i2s-endian-chan-valid-data} 描述了不同通道有效数据位宽下，该寄存器对读取数据的控制。

\begin{table}[H]
    \centering
    \caption{通道有效数据字节序控制}
    \label{tab:i2s-endian-chan-valid-data}
    \begin{tabular}{|p{3.2cm}|R{2.7cm}|R{3.2cm}|R{3.5cm}|}
    \hline
    \rowcolor{lightgray}
    \textbf{通道有效数据位宽} & \textbf{原始数据}  & \textbf{控制后数据字节序}  & \textbf{\hyperref[fielddesc:I2STXBIGENDIAN]{I2S\_TX\_BIG\_ENDIAN}} \\ \hline
    \multirow{2}*{32} & \multirow{2}*{\{B3, B2, B1, B0\}} & \{B3, B2, B1, B0\} & 0 \\\cline{3-4}
                      && \{B0, B1, B2, B3\} & 1 \\\hline

    \multirow{2}*{24} & \multirow{2}*{\{B2, B1, B0\}} & \{B2, B1, B0\} & 0 \\\cline{3-4}
                      && \{B0, B1, B2\} & 1 \\\hline

    \multirow{2}*{16} & \multirow{2}*{\{B1, B0\}} & \{B1, B0\} & 0 \\\cline{3-4}
                      && \{B0, B1\} & 1 \\\hline

                8     & \{B0\} & \{B0\} & x \\\hline
    \end{tabular}
\end{table}

\begin{tiplisting}
B0、B1、B2、B3 各代表一个 8 位数据，符号 \{\} 表示这些字节被组合到一起。例如，\{B3, B2, B1, B0\} 代表一个 32 位的数据，其中 B0 代表 0-7 位，B1 代表 8-15 位，B2 代表 16-23 位，B3 代表 24-31 位。
\end{tiplisting}

\subsubsection{A 律/$\mu$ 律压缩/解压缩}

\chipname{} I2S 按照 32 位（缺省高位补 0$^*$）的方式对排列好字节序的有效数据进行 A 律/$\mu$ 律压缩/解压缩。

\vspace{-1em}
\begin{tiplisting}
缺省高位补 0，即如果有效数据位宽小于 32，则待压缩/解压缩数据的 [31: 有效位宽] 位自动填充 0。
\end{tiplisting}

配置 \hyperref[fielddesc:I2STXPCMBYPASS]{I2S\_TX\_PCM\_BYPASS} 为：
\begin{itemize}
\item 0：进行压缩/解压缩
\item 1：不进行压缩/解压缩
\end{itemize}
配置 \hyperref[fielddesc:I2STXPCMCONF]{I2S\_TX\_PCM\_CONF} 为：
\begin{itemize}
\item 0：A 律解压缩
\item 1：A 律压缩
\item 2：$\mu$ 律解压缩
\item 3：$\mu$ 律压缩
\end{itemize}
至此，数据格式控制的第一阶段完成。

\subsubsection{通道发送数据位宽} 
\label{sec:i2s-tdm-chan-bits}
\chipname{} I2S\regindex{n} 中，\hyperref[fielddesc:I2STXTDMCHANBITS]{I2S\_TX\_TDM\_CHAN\_BITS} 决定了每个通道发送数据的位宽。

\begin{itemize}
\item 当每个通道发送数据的位宽大于有效数据位宽时，会在剩余的位置补充 0。此时，配置 \hyperref[fielddesc:I2STXLEFTALIGN]{I2S\_TX\_LEFT\_ALIGN} 为：
\begin{itemize}
\item 0：有效数据位于发送数据低位，在数据高位补零；
\item 1：有效数据位于发送数据高位，在数据低位补零。
\end{itemize}
\item 当每个通道发送数据的位宽小于有效数据位宽时，每次发送数据仅发送低位有效数据部分，高位部分将被舍弃。
\end{itemize}
至此，数据格式控制的第二阶段完成。

\subsubsection{通道数据比特顺序}
\chipname{} I2S\regindex{n} 通道数据比特顺序由 \hyperref[fielddesc:I2STXBITORDER]{I2S\_TX\_BIT\_ORDER} 控制：
\begin{itemize}
\item 0：不调换有效数据比特顺序；
\item 1：调换有效数据比特顺序。
\end{itemize}
至此，数据格式控制部分全部完成。经过数据格式控制的数据将从高位向低位依次发送。图 \ref{fig:i2s-tx-data-format-control} 所示即为一次完整的 TX 数据格式控制过程。


\begin{figure}[H]
    \centering
    \includegraphics[width=0.8\textwidth]{\modulefiles/figures/i2s_tx_data_config.png}
    \caption{TX 数据格式控制}
    \label{fig:i2s-tx-data-format-control}
\end{figure}


\subsection{通道模式控制}
\chipname{} I2S\regindex{n} 支持 TDM 和 PDM 两种发送模式。置位 \hyperref[fielddesc:I2STXTDMEN]{I2S\_TX\_TDM\_EN} 使能 TDM 输出模式，置位 \hyperref[fielddesc:I2STXPDMEN]{I2S\_TX\_PDM\_EN} 使能 PDM 输出模式。

\begin{tiplisting}
    \hyperref[fielddesc:I2STXTDMEN]{I2S\_TX\_TDM\_EN} 和 \hyperref[fielddesc:I2STXPDMEN]{I2S\_TX\_PDM\_EN} 不能同时置位或同时清零。 
\end{tiplisting}

\subsubsection{TDM 输出模式} \label{sec:i2s-chan-control-tdm-mode}
在 TDM 输出模式下，I2S\regindex{n} 最多支持 16 个发送通道。发送通道数由 \hyperref[fielddesc:I2STXTDMTOTCHANNUM]{I2S\_TX\_TDM\_TOT\_CHAN\_NUM} 控制。例如，配置 \hyperref[fielddesc:I2STXTDMTOTCHANNUM]{I2S\_TX\_TDM\_TOT\_CHAN\_NUM} 为 5，则六个通道（通道 0 \verb+~+ 5）将用于发送数据，见图
\ref{fig:i2s-tdm-chan-control}。 

\vspace{-1em}
\begin{tiplisting}
将 I2S\regindex{n} 模块设置为 TDM 双通道模式，可实现控制大多数 I2S\regindex{n} 双声道编解码器。
\end{tiplisting}

在发送数据的通道中，如果其对应的 \hyperref[fielddesc:I2STXTDMCHAN0EN]{I2S\_TX\_TDM\_CHAN{\textcolor{red}{n}}\_EN} 为：

\begin{itemize}
    \item 1：则该通道发送通道数据；
    \item 0：则该通道发送数据由 \hyperref[fielddesc:I2Ssec:i2s-tx-dataEQUAL]{I2S\_TX\_CHAN\_EQUAL} 控制，当该值为：
    \begin{itemize}
        \item 1：则发送前一个通道的数据；
        \item 0：则发送 \hyperref[fielddesc:I2SSINGLEDATA]{I2S\_SINGLE\_DATA} 的值。
    \end{itemize}
\end{itemize}
当 I2S\regindex{n} 处于主机 TDM 输出模式下，WS 信号由 \hyperref[fielddesc:I2STXWSIDLEPOL]{I2S\_TX\_WS\_IDLE\_POL} 和 \hyperref[fielddesc:I2STXTDMWSWIDTH]{I2S\_TX\_TDM\_WS\_WIDTH} 控制：
    \begin{itemize}
        \item \hyperref[fielddesc:I2STXWSIDLEPOL]{I2S\_TX\_WS\_IDLE\_POL} 的值为 WS 信号的默认电平；
        \item \hyperref[fielddesc:I2STXTDMWSWIDTH]{I2S\_TX\_TDM\_WS\_WIDTH} 的值为在发送所有通道数据的过程中，WS 为默认电平的周期数。
    \end{itemize}

另外，\hyperref[fielddesc:I2STXHALFSAMPLEBITS]{I2S\_TX\_HALF\_SAMPLE\_BITS} 的值乘以 2，即为一个 WS 周期对应的 BCK 周期数。

\textbf{TDM 通道配置示例}

在本示例中，寄存器的配置如下：

\begin{itemize}
    \item 配置 I2S\_TX\_TDM\_CHAN\_NUM 为 5，即选择使用通道 0 \verb+~+ 5 进行数据发送。
    \item 配置 I2S\_TX\_CHAN\_EQUAL 为 1，即如果相应通道的 I2S\_TX\_TDM\_CHAN\regindex{n}\_EN (\regindex{n} = 0 \verb+~+ 5) 被清除，则该通道将发送前一通道的数据。
    \item 配置 I2S\_TX\_TDM\_CHAN0/2/5\_EN 为 1，即这些通道将用于发送通道数据。
    \item 配置 I2S\_TX\_TDM\_CHAN1/3/4\_EN 为 0，则这些通道将发送前一个通道的数据。
\end{itemize}

配置完成后，则数据将按下图方式进行发送。

\begin{figure}[H]
    \centering
    \includegraphics[width=1.0\textwidth]{\modulefiles/figures/i2s_tdm_channel_config.pdf}
    \caption{TDM 通道控制}
    \label{fig:i2s-tdm-chan-control}
\end{figure}


\subsubsection{PDM 输出模式} 

在 PDM 输出模式下，I2S\regindex{n} 可以发送原始 PDM 数据，或者将 PCM 数据转换为 PDM 数据后发送（后者仅 I2S\regindex{0} 支持）。

在 PDM 输出模式下，从 GDMA 取得数据的过程由 \hyperref[fielddesc:I2STXMONO]{I2S\_TX\_MONO} 和 \hyperref[fielddesc:I2STXMONOFSTVLD]{I2S\_TX\_MONO\_FST\_VLD} 控制，具体如下表。请根据内存中存储的数据为单/双通道数据来配置该寄存器。

\begin{table}[H]
    \centering
    \caption{PDM 输出模式下 I2S\regindex{n} 取数逻辑}
    \label{tab:i2s-data-fetch-control-pdm-mode}
    \begin{tabular}{|p{6.6cm}|l|r|r|}
    \hline
    \rowcolor{lightgray}
    \textbf{取数逻辑} &\textbf{模式} &\textbf{\hyperref[fielddesc:I2STXMONO]{I2S\_TX\_MONO}}  & \textbf{\hyperref[fielddesc:I2STXMONOFSTVLD]{I2S\_TX\_MONO\_FST\_VLD}} \\ \hline
                        每个 WS 沿都向 GDMA 发起取数请求 & 双声道 & 0 & x        \\ \hline
                        只在 WS 后半周期向 GDMA 发起取数请求  & 单声道 & 1 & 0   \\ \hline
                        只在 WS 前半周期向 GDMA 发起取数请求 & 单声道 & 1 & 1      \\ \hline
    \end{tabular}
\end{table}

当 I2S\regindex{n} 处于主机 PDM 输出模式时，WS 信号的默认电平由 \hyperref[fielddesc:I2STXWSIDLEPOL]{I2S\_TX\_WS\_IDLE\_POL} 控制，WS 信号频率为 BCK 信号频率的一半。请参考 \ref{sec:i2s-tx-rx-clk} 小节配置 BCK 信号的方式来配置 WS 信号，见图
\ref{fig:i2s-pdm-chan-control-example}。

在\textbf{发送原始 PDM 数据} 时，I2S\regindex{n} 通道数据由 \hyperref[fielddesc:I2STXCHANMOD]{I2S\_TX\_CHAN\_MOD} 和 \hyperref[fielddesc:I2STXWSIDLEPOL]{I2S\_TX\_WS\_IDLE\_POL} 控制，具体如下表。

\begin{table}[H]
    \centering
    \caption{原始 PDM 数据的 I2S\regindex{n} 通道模式控制}
    \label{tab:i2s-chan-control-pdm-mode}
    \begin{threeparttable}
    \begin{tabular}{|p{2.0cm}|L{3.8cm}|L{3.8cm}|R{2.2cm}|R{2.0cm}|}
    \hline
    \rowcolor{lightgray}
    \textbf{模式} & \textbf{左声道} &\textbf{右声道} & \textbf{模式控制字段}\tnote{1} & \textbf{声道选择位}\tnote{2}\\ \hline
    双声道 & 发送左通道数据 & 发送右通道数据 &0 & x \\\hline
                         \multirow{8}*{单声道}  & 发送左通道数据 & 发送左通道数据 & 1 & 0         \\ \cline{2-5}
                                           & 发送右通道数据 & 发送右通道数据 & 1 & 1\\ \cline{2-5}
                                           & 发送右通道数据 & 发送右通道数据 & 2 & 0 \\ \cline{2-5}
                                           & 发送左通道数据 & 发送左通道数据 &2 & 1 \\ \cline{2-5}
                                           & 发送 ``single" 值\tnote{3} & 发送右通道数据 & 3 & 0 \\ \cline{2-5}
                                           & 发送左通道数据 & 发送 ``single" 值 &3 & 1 \\ \cline{2-5}
                                           & 发送左通道数据 & 发送 ``single" 值 &4 & 0 \\ \cline{2-5}
                                           & 发送 ``single" 值 & 发送右通道数据 & 4 & 1 \\ \hline
    \end{tabular}
            \begin{tablenotes}
            \item[1] \hyperref[fielddesc:I2STXCHANMOD]{I2S\_TX\_CHAN\_MOD}
            \item[2] \hyperref[fielddesc:I2STXWSIDLEPOL]{I2S\_TX\_WS\_IDLE\_POL}
            \item[3] single 值等于 \hyperref[fielddesc:I2SSINGLEDATA]{I2S\_SINGLE\_DATA} 的值。
        \end{tablenotes}

    \end{threeparttable}
\end{table}

在 I2S\regindex{0} \textbf{启用 PCM-to-PDM 转换器}时，可以将 GDMA 中的 PCM 数据转为 PDM 数据，并按照 PDM 信号格式输出，配置 \hyperref[fielddesc:I2SPCM2PDMCONVEN]{I2S\_PCM2PDM\_CONV\_EN} 启动该转换器。PCM-to-PDM 转换器的寄存器配置如下：

\begin{itemize}
    \item 配置一线 PDM 输出格式或一/二线 DAC 输出格式，具体如下表。
\begin{table}[H]
    \centering
    \caption{I2S\regindex{0} PCM-to-PDM 数据输出}
    \label{tab:i2s-pcm-pdm-output-mode}
    \begin{tabular}{|p{4.3cm}|R{5.0cm}|R{4.8cm}|}
    \hline
    \rowcolor{lightgray}
    \textbf{通道输出格式}&\textbf{\hyperref[fielddesc:I2STXPDMDACMODEEN]{I2S\_TX\_PDM\_DAC\_MODE\_EN}} & \textbf{\hyperref[fielddesc:I2STXPDMDAC2OUTEN]{I2S\_TX\_PDM\_DAC\_2OUT\_EN}}\\ \hline
    一线 PDM 输出格式\hyperref[i2s-note-1]{$^1$}& 0 & x \\ \hline
    一线 DAC 输出格式\hyperref[i2s-note-2]{$^2$} &1   & 0  \\ \hline 
    二线 DAC 输出格式 & 1 & 1 \\ \hline
    \end{tabular}
\end{table}
\vspace{-1em}
\begin{tiplisting}
\vspace{-2em}
\begin{enumerate}
    \item \label{i2s-note-1} 此处定义的 PDM 输出格式是指一个 WS 周期发送两个通道的 SD 数据。
    \item \label{i2s-note-2} 此处定义的 DAC 输出格式是指一个 WS 周期发送一个通道的 SD 数据。
    \item 一线 PDM/DAC 输出格式使用 I2S\regindex{n}O\_Data\_out 作为数据输出线。
    \item 二线 DAC 输出格式使用 I2S\regindex{n}O\_Data\_out 和 I2S\regindex{n}O\_Data1\_out 作为数据输出线。
    \item 从应用上来说，PDM 输出格式针对的是需要时钟信号、可以解码 PDM 格式编解码器的应用场景，DAC 输出格式针对的是不依赖时钟信号、没有 PDM 格式编解码器、且通过低通滤波可以直接恢复成模拟波形的应用场景。
\end{enumerate}
\end{tiplisting}
\vspace{-1em}

    \item 配置采样频率和上采样率。\newline
    I2S\regindex{0} 启用 PCM-to-PDM 转换器进行数据输出时，PDM 时钟频率即为 BCK 时钟频率。采样频率（$f_{\textrm{Sampling}}$，即 PCM 数据采样频率）与 BCK 时钟频率（PDM 数据采样频率）的关系如下：
    $$f_{\textrm{Sampling}} = \frac{f_{\textrm{BCK}}}{\textrm{OSR}}$$
    
    其中，上采样率 OSR 和 \hyperref[fielddesc:I2STXPDMSINCOSR2]{I2S\_TX\_PDM\_SINC\_OSR2} 的关系为：
    $${\textrm{OSR}} = {\textrm{\hyperref[fielddesc:I2STXPDMSINCOSR2]{I2S\_TX\_PDM\_SINC\_OSR2}}} \times {\textrm{64}}$$
    
    采样频率 $f_{\textrm{Sampling}}$ 和 \hyperref[fielddesc:I2STXPDMFS]{I2S\_TX\_PDM\_FS} 的对应关系为：
    $$f_{\textrm{Sampling}} = {\textrm{\hyperref[fielddesc:I2STXPDMFS]{I2S\_TX\_PDM\_FS}}} \times {\textrm{100}}$$
    
    请根据需要的采样频率、上采样率以及 PDM 时钟频率进行寄存器配置。
\end{itemize}


\textbf{PDM 通道配置示例}

本示例针对 I2S\regindex{0}，示例中寄存器的配置如下：

\begin{itemize}
    \item 配置 \hyperref[fielddesc:I2SPCM2PDMCONVEN]{I2S\_PCM2PDM\_CONV\_EN} 为 0，即选择输出原始 PDM 数据。
    \item 配置 \hyperref[fielddesc:I2STXMONO]{I2S\_TX\_MONO} 为 0，即 I2S\regindex{0} 在 WS 的高低电平均通过 GDMA 从存储器取数。
    \item 配置 \hyperref[fielddesc:I2STXCHANMOD]{I2S\_TX\_CHAN\_MOD} 为 2，即选择单声道模式。
    \item 配置 \hyperref[fielddesc:I2STXWSIDLEPOL]{I2S\_TX\_WS\_IDLE\_POL} 为 1，则左声道和右声道均发送左通道数据，右声道数据将被舍弃。
\end{itemize}

配置完成后，假定存储器中的数据\textbf{经过数据格式控制之后}为：

\begin{table}[H]
    \centering
    \begin{tabular}{|c|c|c|c|c|c|c|}
    \hline
    Left & Right & Left & Right & ... & Left & Right \\ \hline
    \end{tabular}
\end{table}
\vspace{-0.5em}
\begin{tiplisting}
\vspace{-2em}
\begin{enumerate}
    \item 此处的数据并非存储器中的原始数据，而是经过数据格式控制之后的数据。
    \item 表中的 ``Left" 和 ``Right" 代表通道数据，其位宽为通道有效数据位宽，请参考章节 \ref{sec:i2s-data-control}。
\end{enumerate}
\end{tiplisting}
\vspace{-1em}
则数据经过通道模式控制后，将按照下图方式进行发送。

\begin{figure}[H]
    \centering
    \includegraphics[width=0.8\textwidth]{\modulefiles/figures/i2s_chan_mode_20220210.pdf}
    \caption{PDM 通道控制}
    \label{fig:i2s-pdm-chan-control-example}
\end{figure}

\subsection{同步计数器}

\chipname{} I2S\regindex{n} TX 模块包含同步计数器，用于反映当前 I2S\regindex{n} 的数据发送情况。同步计数器可用于检测 I2S\regindex{n} 发送数据错误，或在多个设备使用 I2S\regindex{n} 发送数据时进行同步。

\chipname{} I2S\regindex{n} 共包含两种类型的同步计数器：

\begin{itemize}
    \item \hyperref[fielddesc:I2STXFIFOCNT]{I2S\_TX\_FIFO\_CNT}
    \begin{itemize}
        \item 每次 I2S\regindex{n} TX FIFO 成功读取一个数据，该计数器加 1。
        \item 置位 \hyperref[fielddesc:I2STXFIFOCNTRST]{I2S\_TX\_FIFO\_CNT\_RST} 可复位该同步计数器。
        \item 该计数器位宽为 31 比特，溢出时自动归零。
    \end{itemize}
    \item \hyperref[fielddesc:I2STXBCKCNT]{I2S\_TX\_BCK\_CNT}
    \begin{itemize}
        \item 每次 I2S\regindex{n} TX 发送一个 BCK 周期，该计数器加 1。
        \item 可以通过置位 \hyperref[fielddesc:I2STXBCKCNTRST]{I2S\_TX\_BCK\_CNT\_RST} 来复位该同步计数器。
        \item 该计数器位宽为 31 比特，溢出时自动归零。
    \end{itemize}
\end{itemize}

\vspace{-0.5em}
\begin{tiplisting}
\vspace{-2em}
\begin{itemize}
    \item 每次 I2S\regindex{n} TX FIFO 发送一个数据，均意味着发送了一个通道数据，因此，当 I2S\regindex{n} 正确发送数据时，\hyperref[fielddesc:I2STXBCKCNT]{I2S\_TX\_BCK\_CNT} = \hyperref[fielddesc:I2STXFIFOCNT]{I2S\_TX\_FIFO\_CNT} × 通道发送数据位宽（请参考小节 \ref{sec:i2s-tdm-chan-bits}）。
    \item 由于 I2S\regindex{n} TX 模块发送数据时，从 FIFO 取出数据到把数据发送至线上存在一些延迟，因此读到 \hyperref[fielddesc:I2STXBCKCNT]{I2S\_TX\_BCK\_CNT} 的值与 \hyperref[fielddesc:I2STXFIFOCNT]{I2S\_TX\_FIFO\_CNT} 的值可能略微不对应。
\end{itemize}
\end{tiplisting}
\vspace{-1em}

\textbf{应用场景举例}
\begin{enumerate}
    \item 多个设备 I2S TX 同步\\
    多个设备同时通过 I2S TX 发送数据时，由于时钟源的微小频率差，长时间发送数据后会出现同步问题。此时，可以通过读取每个设备的 \hyperref[fielddesc:I2STXFIFOCNT]{I2S\_TX\_FIFO\_CNT} 值来判断发送音频数据的位置，从而进行同步操作。
    \item I2S TX 稳定性检查\\
    当设备处于 GDMA 占用较多的应用场景时，可能会由于多路仲裁，导致 I2S TX 出现偶尔缺数。此时，检查 \hyperref[fielddesc:I2STXBCKCNT]{I2S\_TX\_BCK\_CNT} 与 \hyperref[fielddesc:I2STXFIFOCNT]{I2S\_TX\_FIFO\_CNT} 的比例关系，可以确认设备运行 I2S TX 时的稳定性。
\end{enumerate}


\section{接收数据}\label{sec:i2s-rx-data}

I2S\regindex{n} 接收数据时，从外设接口读取数据，经过通道模式控制和数据格式控制，通过 GDMA 输入数据至内存。

\subsection{通道模式控制}

\chipname{} I2S\regindex{n} 支持 TDM 和 PDM 两种输入模式。置位 \hyperref[fielddesc:I2SRXTDMEN]{I2S\_RX\_TDM\_EN} 使能 TDM 输入模式，置位 \hyperref[fielddesc:I2SRXPDMEN]{I2S\_RX\_PDM\_EN} 使能 PDM 输入模式。

\begin{tiplisting}
\hyperref[fielddesc:I2SRXTDMEN]{I2S\_RX\_TDM\_EN} 和 \hyperref[fielddesc:I2SRXPDMEN]{I2S\_RX\_PDM\_EN} 不能同时置位或同时清零。
\end{tiplisting}
\vspace{-1em}

\subsubsection{TDM 输入模式}

在 TDM 输入模式下，I2S\regindex{n} 最多支持 16 个接收通道。接收通道数由 \hyperref[fielddesc:I2SRXTDMTOTCHANNUM]{I2S\_RX\_TDM\_TOT\_CHAN\_NUM} 控制。例如，配置 \hyperref[fielddesc:I2SRXTDMTOTCHANNUM]{I2S\_RX\_TDM\_TOT\_CHAN\_NUM} 为 5，则通道 0 \verb+~+ 5 接收数据。

在接收数据的通道中，如果其对应的 \hyperref[fielddesc:I2STXTDMCHAN0EN]{I2S\_RX\_TDM\_CHAN{\textcolor{red}{n}}\_EN} 为：
\begin{itemize}
    \item 0：则该通道数据无效，不会存入 RX FIFO；
    \item 1：则该通道数据有效，会被存入 RX FIFO。
\end{itemize}
当 I2S\regindex{n} 处于主机 TDM 输入模式下，WS 信号由 \hyperref[fielddesc:I2SRXWSIDLEPOL]{I2S\_RX\_WS\_IDLE\_POL} 和 \hyperref[fielddesc:I2SRXTDMWSWIDTH]{I2S\_RX\_TDM\_WS\_WIDTH} 控制：
\begin{itemize}
    \item \hyperref[fielddesc:I2SRXWSIDLEPOL]{I2S\_RX\_WS\_IDLE\_POL} 的值为 WS 信号的默认电平；
    \item \hyperref[fielddesc:I2SRXTDMWSWIDTH]{I2S\_RX\_TDM\_WS\_WIDTH} 的值为在接收所有通道的过程中，WS 为默认电平的周期数。
\end{itemize}
另外，\hyperref[fielddesc:I2SRXHALFSAMPLEBITS]{I2S\_RX\_HALF\_SAMPLE\_BITS} 的值乘以 2，即为一个 WS 周期对应的 BCK 周期数。

\subsubsection{PDM 输入模式} 

在 PDM 输入模式下，I2S\regindex{n} 可以接收原始 PDM 数据，或者将 PDM 数据转换为 PCM 数据后接收（后者仅 I2S\regindex{0} 支持）。

在 PDM 输入模式下，I2S\regindex{n} 将通道中的串行数据转换为待输入数据。

当 I2S\regindex{n} 处于主机 PDM 输入模式下，WS 信号的默认电平由 \hyperref[fielddesc:I2SRXWSIDLEPOL]{I2S\_RX\_WS\_IDLE\_POL} 控制，WS 信号频率为 BCK 信号频率的一半，请参考 \ref{sec:i2s-tx-rx-clk} 小节配置 BCK 信号的方式配置 WS 信号。注意，在 PDM 输入模式下，请配置 \hyperref[fielddesc:I2SRXHALFSAMPLEBITS]{I2S\_RX\_HALF\_SAMPLE\_BITS} 的值与 \hyperref[fielddesc:I2SRXBITSMOD]{I2S\_RX\_BITS\_MOD} 的值相同。

在 I2S\regindex{0} \textbf{启用 PDM-to-PCM 转换器}时，可以将接收的 PDM 数据转为 PCM 数据，再进行数据控制。配置 \hyperref[fielddesc:I2SRXPDM2PCMEN]{I2S\regindex{0}\_RX\_PDM2PCM\_EN} 启用该转换器。PDM-to-PCM 转换器的寄存器配置如下：

\begin{itemize}
    \item 配置采样频率和下采样率。\\
    I2S\regindex{0} 使用 PDM-to-PCM 转换器接收数据时，PDM 时钟频率与主从机模式的关系如下：
    \begin{itemize}
        \item 主机模式：PDM 时钟频率即为 BCK 时钟频率。
        \item 从机模式：PDM 时钟频率由外部提供。
    \end{itemize}
    采样频率 ($f_{\textrm{Sampling}}$) 与 PDM 时钟频率的关系如下：
    $$f_{\textrm{Sampling}} = \frac{f_{\textrm{PDM}}}{\textrm{DSR}}$$
    其中，下采样率 DSR 和 \hyperref[fielddesc:I2SRXPDMSINCDSR16EN]{I2S\_RX\_PDM\_SINC\_DSR\_16\_EN} 的关系为：
    $${\textrm{DSR}} = {\textrm{\hyperref[fielddesc:I2SRXPDMSINCDSR16EN]{I2S\_RX\_PDM\_SINC\_DSR\_16\_EN}}} \times {\textrm{64}}$$
    请根据需要的主从机模式、采样频率以及下采样率进行寄存器配置。
    \item 配置有效通道。\\
    启用 PDM-to-PCM 转换器时，最多支持 8 个通道的输入信号，寄存器配置以及对应通道如下表：
\begin{table}[H]
    \centering
    \caption{PDM-to-PCM 数据输入}
    \label{table:I2S_PDM_RX_MODE}
    \begin{tabular}{|p{4.3cm}|p{2.0cm}|p{5.8cm}|}
    \hline
    \rowcolor{lightgray}
    \textbf{输入数据信号}&\textbf{通道} &\textbf{使能寄存器} \\ \hline
    \multirow{2}*{I2S\regindex{0}I\_Data\_in}  & 左声道 & \hyperref[fielddesc:I2SRXTDMPDMCHAN0EN]{I2S\_RX\_TDM\_PDM\_CHAN0\_EN} \\ \cline{2-3}
                                               & 右声道 & \hyperref[fielddesc:I2SRXTDMPDMCHAN0EN]{I2S\_RX\_TDM\_PDM\_CHAN1\_EN} \\ \hline
    \multirow{2}*{I2S\regindex{0}I1\_Data\_in} & 左声道 & \hyperref[fielddesc:I2SRXTDMPDMCHAN0EN]{I2S\_RX\_TDM\_PDM\_CHAN2\_EN} \\ \cline{2-3}
                                               & 右声道 & \hyperref[fielddesc:I2SRXTDMPDMCHAN0EN]{I2S\_RX\_TDM\_PDM\_CHAN3\_EN} \\ \hline
    \multirow{2}*{I2S\regindex{0}I2\_Data\_in} & 左声道 & \hyperref[fielddesc:I2SRXTDMPDMCHAN0EN]{I2S\_RX\_TDM\_PDM\_CHAN4\_EN} \\ \cline{2-3}
                                               & 右声道 & \hyperref[fielddesc:I2SRXTDMPDMCHAN0EN]{I2S\_RX\_TDM\_PDM\_CHAN5\_EN} \\ \hline
    \multirow{2}*{I2S\regindex{0}I3\_Data\_in} & 左声道 & \hyperref[fielddesc:I2SRXTDMPDMCHAN0EN]{I2S\_RX\_TDM\_PDM\_CHAN6\_EN} \\ \cline{2-3}
                                               & 右声道 & \hyperref[fielddesc:I2SRXTDMPDMCHAN0EN]{I2S\_RX\_TDM\_PDM\_CHAN7\_EN} \\ \hline
    \end{tabular}
\end{table} 
\end{itemize}

\subsection{数据格式控制}

数据格式控制分为两个阶段：
\begin{itemize}
    \item 第一阶段：从串行数据流输入转换为待输入数据，并存入 RX FIFO；
    \item 第二阶段：将待输入数据从 RX FIFO 中读出，并进行输入数据模式转换。
\end{itemize}

\subsubsection{通道数据比特顺序} 

\chipname{} I2S\regindex{n} 通道数据将从高位到低位依次存入待输入数据，并由 \hyperref[fielddesc:I2SRXBITORDER]{I2S\_RX\_BIT\_ORDER} 控制是否调换比特序，当该值为：
\begin{itemize}
\item 0：不调换待输入数据的比特序；
\item 1：调换待输入数据的比特序。
\end{itemize}
至此，数据格式控制的第一阶段完成，经过比特顺序控制的待输入数据被存入 RX FIFO。

\subsubsection{通道储存数据位宽}

\hyperref[fielddesc:I2SRXBITSMOD]{I2S\_RX\_BITS\_MOD} 和 \hyperref[fielddesc:I2SRX24FILLEN]{I2S\_RX\_24\_FILL\_EN} 决定了每个通道的储存数据位宽，其可取值和对应的储存数据位宽如下表：

\begin{table}[H]
    \centering
    \caption{通道储存数据位宽控制}
    \label{tab:i2s-chan-storage-data-width}
    \begin{tabular}{|p{3cm}|R{3.3cm}|R{4cm}|}
    \hline
    \rowcolor{lightgray}
    \textbf{通道储存数据位宽} & \textbf{\hyperref[fielddesc:I2SRXBITSMOD]{I2S\_RX\_BITS\_MOD}} & \textbf{\hyperref[fielddesc:I2SRX24FILLEN]{I2S\_RX\_24\_FILL\_EN}}  \\ \hline
    \multirow{2}*{32} & 31 & x \\\cline{2-3}
                 & 23 & 1 \\ \hline
    24 & 23 & 0 \\\hline
    16    & 15 & x \\\hline
    8     & 7 & x \\\hline
    \end{tabular}
\end{table}

\subsubsection{通道接收数据位宽} 

\hyperref[fielddesc:I2SRXTDMCHANBITS]{I2S\_RX\_TDM\_CHAN\_BITS} 决定了每个通道接收数据的位宽。

\begin{itemize}
\item 当每个通道储存数据位宽小于接收数据的位宽时，接收时仅会在接收数据中取储存位宽作为储存数据。此时，配置 \hyperref[fielddesc:I2SRXLEFTALIGN]{I2S\_RX\_LEFT\_ALIGN} 为：
\begin{itemize}
\item 0：储存数据位于接收数据低位；
\item 1：储存数据位于接收数据高位。
\end{itemize}
\item 当每个通道接收数据的位宽小于储存数据位宽时，会将接收数据高位补 0，作为储存数据。
\end{itemize}

\subsubsection{通道储存数据字节序} 

通道接收数据经过截取/补全后，成为了待储存数据，\hyperref[fielddesc:I2SRXBIGENDIAN]{I2S\_RX\_BIG\_ENDIAN} 用于控制待储存数据的字节序。下表描述了不同通道储存数据位宽下，该寄存器对待储存数据的控制。
\begin{table}[H]
    \centering
    \caption{通道储存数据字节序控制}
    \label{tab:i2s-chan-storage-data-endian}
    \begin{tabular}{|p{3cm}|R{3cm}|R{3.5cm}|R{3.5cm}|}
    \hline
    \rowcolor{lightgray}
        \textbf{通道储存数据位宽} & \textbf{原始数据} & \textbf{控制后数据} & \textbf{\hyperref[fielddesc:I2SRXBIGENDIAN]{I2S\_RX\_BIG\_ENDIAN}} \\ \hline
 \multirow{2}*{32} & \multirow{2}*{\{B3, B2, B1, B0\}} & \{B3, B2, B1, B0\} & 0 \\\cline{3-4}
                      && \{B0, B1, B2, B3\} & 1 \\\hline

    \multirow{2}*{24} & \multirow{2}*{\{B2, B1, B0\}} & \{B2, B1, B0\} & 0 \\\cline{3-4}
                      && \{B0, B1, B2\} & 1 \\\hline

    \multirow{2}*{16} & \multirow{2}*{\{B1, B0\}} & \{B1, B0\} & 0 \\\cline{3-4}
                      && \{B0, B1\} & 1 \\\hline

                8     & \{B0\} & \{B0\} & x \\\hline
    \end{tabular}
\end{table}

\subsubsection{A 律/$\mu$ 律压缩/解压缩} 

I2S\regindex{n} 对排列好字节序的待储存数据会按照 32 位（缺省高位补 0）的方式进行 A 律/$\mu$ 律压缩/解压缩。

配置 \hyperref[fielddesc:I2SRXPCMBYPASS]{I2S\_RX\_PCM\_BYPASS} 为：
\begin{itemize}
\item 0：进行压缩/解压缩
\item 1：不进行压缩/解压缩
\end{itemize}
配置 \hyperref[fielddesc:I2SRXPCMCONF]{I2S\_RX\_PCM\_CONF} 为：
\begin{itemize}
\item 0：A 律解压缩
\item 1：A 律压缩
\item 2：$\mu$ 律解压缩
\item 3：$\mu$ 律压缩
\end{itemize}

至此，数据格式控制部分全部完成，数据通过 GDMA 存入内存。

\section{事件任务矩阵功能} \label{sec:i2s-etm-func}

在 \chipname{} 中，I2S\regindex{n} 支持 ETM 功能，即可以通过任意外设的 ETM 事件触发 I2S\regindex{n} 的 ETM 任务，或者通过 I2S\regindex{n} 的 ETM 事件触发任意外设的 ETM 任务。关于 ETM 更多详细信息，请参考章节 \ref{mod:evntaskmatrix} \nameref{mod:evntaskmatrix}。这里仅介绍与 I2S\regindex{n} 相关的 ETM 任务和 ETM 事件。

I2S\regindex{n} 可接收的 ETM 任务有：

\begin{itemize}
    \item \label{task:I2STASKSTARTTX}I2S\regindex{n}\_TASK\_START\_TX：触发时启动 I2S\regindex{n} TX 数据传输。
    \item \label{task:I2STASKSTARTRX}I2S\regindex{n}\_TASK\_START\_RX：触发时启动 I2S\regindex{n} RX 数据传输。
    \item \label{task:I2STASKSTOPTX}I2S\regindex{n}\_TASK\_STOP\_TX：触发时停止 I2S\regindex{n} TX 数据传输。
    \item \label{task:I2STASKSTOPRX}I2S\regindex{n}\_TASK\_STOP\_RX：触发时停止 I2S\regindex{n} RX 数据传输。
\end{itemize}

I2S\regindex{n} 可产生的 ETM 事件有：

\begin{itemize}
    \item \label{evt:I2SEVTTXDONE}I2S\regindex{n}\_EVT\_TX\_DONE：表示 I2S\regindex{n} TX 发送完成。
    \item \label{evt:I2SEVTRXDONE}I2S\regindex{n}\_EVT\_RX\_DONE：表示 I2S\regindex{n} RX 接收完成。
    \item \label{evt:I2SEVTXWORDSSENT}I2S\regindex{n}\_EVT\_X\_WORDS\_SENT：表示 I2S\regindex{n} TX 发送字数大于等于 \hyperref[fielddesc:I2SETMTXSENDWORDNUM]{I2S\_ETM\_TX\_SEND\_WORD\_NUM} 的值。
    \item \label{evt:I2SEVTXWORDSRECEIVED}I2S\regindex{n}\_EVT\_X\_WORDS\_RECEIVED：表示 I2S\regindex{n} RX 接收字数大于等于 \hyperref[fielddesc:I2SETMRXRECEIVEWORDNUM]{I2S\_ETM\_RX\_RECEIVE\_WORD\_NUM} 的值。
\end{itemize}

在具体应用中，可以配置 I2S\regindex{n} 的 ETM 事件来触发 I2S\regindex{n} 的 ETM 任务，例如配置 I2S\regindex{n}\_EVT\_X\_WORDS\_SENT 事件触发 I2S\regindex{n}\_TASK\_STOP\_TX 任务，从而以 ETM 的方式完成停止 I2S\regindex{n} 的操作。


%%%%%%%%%%%%%%%%%%%%%
%%% Interrupts %%%
%%%%%%%%%%%%%%%%%%%%%

\section{I2S 中断} \label{sec:i2s-interrupts}

\chipname{} 的 I2S\regindex{n} 可以生成 I2S\regindex{n}\_INTR 中断信号并发送给\hyperref[mod:intmtrx]{中断矩阵}。这一中断信号由以下 I2S\regindex{n} 的内部中断源生成。

\begin{itemize}
\item \label{int:I2STXHUNGINT}I2S\regindex{n}\_TX\_HUNG\_INT：当发送数据超时，触发此中断。例如，I2S\regindex{n} 配置为从机发送模式，但主机长时间未提供 BCK 或 WS 信号，则将触发该中断。超时配置见寄存器 \hyperref[regdesc:I2SLCHUNGCONFREG]{I2S\_LC\_HUNG\_CONF\_REG}。
\item \label{int:I2SRXHUNGINT}I2S\regindex{n}\_RX\_HUNG\_INT：当接收数据超时，触发此中断。例如，I2S\regindex{n} 配置为从机接收模式，但主机长时间未发送数据，则将触发该中断。超时配置见寄存器 \hyperref[regdesc:I2SLCHUNGCONFREG]{I2S\_LC\_HUNG\_CONF\_REG}。
\item \label{int:I2STXDONEINT}I2S\regindex{n}\_TX\_DONE\_INT：当发送数据完成，触发此中断。
\item \label{int:I2SRXDONEINT}I2S\regindex{n}\_RX\_DONE\_INT：当接收数据完成，触发此中断。
\end{itemize}

\begin{tiplisting}
中断、中断源、中断信号等术语的解释以及它们之间的联系请参考章节 \ref{mod:intmtrx} \textit{\nameref{mod:intmtrx}} > \ref{sec:intmtx-int-terminology} \textit{\nameref{sec:intmtx-int-terminology}}。
\end{tiplisting}

每个中断源都有一组通用的配置寄存器，详见章节 \hyperref[interrupt-config-registers]{\textit{中断配置寄存器}}。具体寄存器名称请查看 \ref{sec:i2s-reg-summ} \textit{\nameref{sec:i2s-reg-summ}}。

\section{软件配置流程}
\subsection{软件配置 I2S 发送流程}\label{sec:i2s-configure-i2s-tx-mode}

软件配置 I2S\regindex{n} 发送的流程如下：

\begin{enumerate}
    \item 根据 \ref{sec:i2s-tx-rx-clk} 小节的描述，配置时钟。
    \item 根据表 \ref{tab:i2s-signal-descr} 的描述，配置信号管脚。
    \item 根据主从机模式配置 \hyperref[fielddesc:I2STXSLAVEMOD]{I2S\_TX\_SLAVE\_MOD}：
    \begin{itemize}
        \item 0：主机发送模式
        \item 1：从机发送模式
    \end{itemize}
    \item 根据 \ref{sec:i2s-tx-data} 小节的描述，配置正确的发送数据模式和发送通道模式，置位 \hyperref[fielddesc:I2STXUPDATE]{I2S\_TX\_UPDATE}。
    \item 根据 \ref{sec:i2s-reset} 小节的描述，复位 TX 模块和发送 FIFO。
    \item 根据 \ref{sec:i2s-interrupts} 小节的描述使能相应的中断。\label{other:i2s-item-tx-enable-corresponding-interrupts}
    \item 配置 GDMA 发送链表。 
    \item 根据需要置位 \hyperref[fielddesc:I2STXSTOPEN]{I2S\_TX\_STOP\_EN}，更多信息见章节 \ref{sec:i2s-master-slave-tx-mode}。
    \item 开始发送数据：
    \begin{itemize}
        \item 主机模式下，等待 I2S\regindex{n} 从设备配置完成后，置位 \hyperref[fielddesc:I2STXSTART]{I2S\_TX\_START} 开始发送数据；
        \item 从机模式下，置位 \hyperref[fielddesc:I2STXSTART]{I2S\_TX\_START}。I2S\regindex{n} 主设备提供 BCK 和 WS 信号后，开始发送数据。
    \end{itemize}
    \item 等待步骤 \ref{other:i2s-item-tx-enable-corresponding-interrupts} 设置的中断信号，或查询 \hyperref[fielddesc:I2STXIDLE]{I2S\_TX\_IDLE} 检查传输是否结束：
    \begin{itemize}
        \item 0：发送设备为工作状态；
        \item 1：发送设备为空闲状态。
    \end{itemize}
    \item 清零 \hyperref[fielddesc:I2STXSTART]{I2S\_TX\_START} 停止发动数据。
\end{enumerate}

\subsection{软件配置 I2S 接收流程}\label{sec:i2s-configure-i2s-rx-mode}

软件配置 I2S\regindex{n} 接收模式的流程如下：

\begin{enumerate}
    \item 根据 \ref{sec:i2s-tx-rx-clk} 小节的描述，配置时钟。
    \item 根据表 \ref{tab:i2s-signal-descr} 的描述，配置信号管脚。
    \item 配置 \hyperref[fielddesc:I2SRXSLAVEMOD]{I2S\_RX\_SLAVE\_MOD} 选择需要的模式：
    \begin{itemize}
        \item 0：主机接收模式
        \item 1：从机接收模式
    \end{itemize}
    \item 根据 \ref{sec:i2s-rx-data} 小节的描述，配置正确的接收通道模式和接收数据模式，置位 \hyperref[fielddesc:I2SRXUPDATE]{I2S\_RX\_UPDATE}。
    \item 根据 \ref{sec:i2s-reset} 小节的描述，复位 RX 模块和接收 FIFO。
    \item 根据 \ref{sec:i2s-interrupts} 小节的描述使能相应的中断。\label{other:i2s-item-rx-enable-corresponding-interrupts}
    \item 配置 GDMA 接收链表，并在 \hyperref[regdesc:I2SRXEOFNUMREG]{I2S\_RX\_EOF\_NUM\_REG} 中配置接收数据长度。
    \item 开始接收数据：
    \begin{itemize}
        \item 在主机模式下，等待从机准备好后，置位 \hyperref[fielddesc:I2SRXSTART]{I2S\_RX\_START} 开始接收数据；
        \item 在从机模式下，置位 \hyperref[fielddesc:I2SRXSTART]{I2S\_RX\_START}，等待主机提供 BCK 和 WS 信号后开始接收数据。
    \end{itemize}
    \item 接收的数据由 GDMA 根据配置存到存储器的指定地址。最终产生步骤 \ref{other:i2s-item-rx-enable-corresponding-interrupts} 中设置的中断。
\end{enumerate}


% >>> If your module has no registers, delete sample text
%       for Sections Register Summary and Registers below <<<

%%%%%%%%%%%%%%%%%%%%%%%%
%%% Register Summary %%%
%%%%%%%%%%%%%%%%%%%%%%%%

% >>> Update the sample text below and insert
%     the info on register summary into the subfile <<<

\newpage
\hypertarget{i2s-reg-summ}{}
\section{寄存器列表}
\label{sec:i2s-reg-summ}

本小节的所有地址均为相对于 I2S\regindex{n} 基地址的地址偏移量（相对地址），具体基地址请见章节 \ref{mod:sysmem} \textit{\nameref{mod:sysmem}} 中的表 \ref{tab:sysmem-base-address}。

请查看章节 \hyperref[glossary-access-types]{\textit{寄存器的访问类型}}，了解“访问”列缩写的含义。

\subfile{\modulefiles/48-I2S-reg-summ__CN}


%%%%%%%%%%%%%%%%%
%%% Registers %%%
%%%%%%%%%%%%%%%%%

% >>> Update the sample text below and insert
%      the info on registers into the subfile <<<

\newpage
\section{寄存器}

本小节的所有地址均为相对于 I2S\regindex{n} 基地址的地址偏移量（相对地址），具体基地址请见章节 \ref{mod:sysmem} \textit{\nameref{mod:sysmem}} 中的表 \ref{tab:sysmem-base-address}。

关于保留 (reserved) 域的处理，请查看章节 \hyperref[programming-reserved-field]{\textit{如何配置寄存器的保留域}}。

\subfile{\modulefiles/48-I2S-reg__CN}

\end{document}
